% =============================================================================
% SECTION <5>: LIFE PROFILE DESCRIPTION
% =============================================================================
\subsection{5. Life Profile Description}

This section defines the chronological lifecycle phases, constellation operational modes, contingency recovery states, and end-of-life disposal profile of Configuration Item \textbf{\texttt{CI-03}: Space Segment (\abbr{SSEG}, \abbr{WBS} Element \texttt{123456.03})} from on-ground integration through post-mission passivation. Detailed lifecycle, outgassing, and disposal trade-off justifications are documented in the Space Segment Design Justification File (\texttt{EWAS-HSB-SSEG-DDJF-03}, \textbf{RD-02}).

\subsubsection{5.1 On-Ground AIT, Storage, and Transport Phase}

In accordance with the Verification Plan (\textbf{AD-02}) and \abbr{WBS} \texttt{123456.04.02}, Space Segment manufacturing and verification follow a phased model philosophy comprising Engineering and Structural-Thermal Models (\texttt{EM}/\texttt{STM}) for Hardware-in-the-Loop (\abbr{HIL}) and structural-thermal qualification (\texttt{R-039}), lead Qualification/Proto-Flight Models (\texttt{QM}/\texttt{PFM}) for full environmental qualification, and recurring Flight Models (\texttt{FM}) organized as 8 plane-lots of 9 \texttt{SATL} units plus 6 individual \texttt{SATG} units. Cryogenic optical payloads and 1,550~nm \abbr{OISL} assemblies are integrated under ISO Class~5 dry gaseous nitrogen ($\mathrm{GN_2}$) purged cleanroom conditions ($\mathrm{CVCM} < 0.01\%$, $\mathrm{TML} < 1.0\%$), while spacecraft bus integration uses ISO Class~8 facilities; spare flight units may be stored up to 24~months and transported in hermetically sealed, climate-controlled containers limiting dynamic transport loads to $\le 3.0\,g$.

\subsubsection{5.2 Launch, LEOP, and In-Orbit Commissioning Phase}

The 72 \texttt{SATL} spacecraft are deployed via 8 dedicated plane-batch launches of 9 spacecraft each mounted on cylindrical rideshare dispenser stacks ($\le 2{,}800.0\text{ kg}$ per stack) into the 1,000~km, $80^\circ$ inclination planes, while the 6 \texttt{SATG} spacecraft launch on single- or dual-manifest heavy-lift vehicles toward 35,786~km equatorial \abbr{GEO}; all hazardous systems remain under at least three independent mechanical or electrical safety inhibits throughout ascent. Following staggered dispenser ejection, each spacecraft autonomously damps post-separation bus tip-off rates from $\le 2.0\text{--}3.0^\circ/\text{s}$ down to $< 0.1^\circ/\text{s}$, deploys solar arrays and antennas, establishes encrypted S-band \abbr{TT\&C} lock, and executes a 72--96~h outgassing decontamination bake-out ($310\text{--}320\text{ K}$) prior to cooling its infrared focal planes to $65\text{--}95\text{ K}$. Each vehicle then executes orbit insertion and phasing burns into its assigned Walker Delta $72/8/1$ ($\Delta u_{\text{inter}} = 22.5^\circ$) or $60^\circ$ \abbr{GEO} slot, completes star-tracker-to-payload boresight alignment and radiometric calibration, locks the intra-/cross-plane and \abbr{GEO}--\abbr{LEO} \abbr{OISL} network, and passes end-to-end validation trials (\texttt{R-040}) before operational handover.

\subsubsection{5.3 Nominal Operational Modes (\texttt{N1}--\texttt{N5})}

Across its $\ge 7\text{-year}$ (\texttt{SATL}) and $\ge 15\text{-year}$ (\texttt{SATG}) nominal design lifetime (\texttt{R-016}), the Space Segment sustains $> 99.9\%$ system availability (\texttt{R-015}) through five segment-level operational modes:
\begin{enumerate}
    \item \textbf{Mode N1 --- Continuous Boost-Phase Surveillance:} Default 24/7 staring mode of all 6 \texttt{SATG} nodes (supplemented by un-cued \texttt{SATL} nodes at $> 60^\circ\text{N}$) executing $\ge 10.0\text{ Hz}$ \abbr{SWIR}/\abbr{MWIR} detection, onboard proton despiking, clutter filtering (\texttt{R-019}, \texttt{R-022}), and \emph{GEO Event Report} downlink within $\le 10.0\text{ s}$ (\texttt{R-010}, \texttt{R-042}).
    \item \textbf{Mode N2 --- Cued Agile Stereoscopic Tracking:} Preemptive tracking mode executed by two Ground-selected \texttt{SATL} nodes upon receiving a \emph{LEO Tracking Request} (\texttt{R-043}, \texttt{R-044}); encompasses $< 5.0\text{ s}$ independent payload \abbr{LOS} lock-on with $< 2.0\text{ arcsec/s}$ stability (\texttt{R-024}), $10\text{--}20\text{ Hz}$ \abbr{OISL} vector exchange for onboard preliminary 3D trajectory estimation (\texttt{R-023}), and dual-product downlink within $\le 30.0\text{ s}$ (\texttt{R-011}, \texttt{R-012}, \texttt{R-045}).
    \item \textbf{Mode N3 --- Mandatory \abbr{OISL} Mesh Relay:} Continuous 24/7 background optical network operation across all 78 nodes maintaining intra-/cross-plane \abbr{LEO} links and \abbr{GEO}--\abbr{LEO} trunks, Priority-0 \abbr{QoS} warning-chain routing, \abbr{COMSEC} key rotation (\texttt{R-031}), and anti-jamming resilience (\texttt{R-032}, \texttt{R-033}).
    \item \textbf{Mode N4 --- Staggered Station-Keeping and Collision Avoidance:} Periodic slot maintenance ($\pm 0.5^\circ$ phase in \abbr{LEO}; $\pm 0.05^\circ$ box in \abbr{GEO}), momentum dumping, and Collision Avoidance Manoeuvres (\abbr{CAM}) triggered by Conjunction Data Messages ($P_c > 10^{-4}$, \texttt{IF-EXT-04}), staggered so at most one \texttt{SATL} node per coverage zone manoeuvres at any time (\texttt{R-008}, \texttt{R-015}).
    \item \textbf{Mode N5 --- In-Orbit Calibration and \abbr{OTA} Software Maintenance:} Scheduled blackbody/deep-space radiometric calibration, thermo-elastic boresight updates, and authenticated over-the-air uplink of archive-refined detection/tracking algorithms (\texttt{R-004}, \texttt{R-036}, \texttt{R-041}).
\end{enumerate}

\subsubsection{5.4 Contingency / Safe Mode and End-of-Life Disposal Phase}

To prevent single-point failures from degrading constellation availability (\texttt{R-015}, \texttt{R-017}), each spacecraft implements a three-level autonomous \abbr{FDIR} architecture: Level~1 provides sub-100~ms transparent recovery from transient radiation upsets (\abbr{EDAC}, \abbr{TMR}, proton despiking, and \abbr{SEL} rail resets); Level~2 executes autonomous redundant unit failover within $\le 2.0\text{ s}$ and broadcasts updated availability across the \abbr{OISL} mesh; and Level~3 transitions faulted spacecraft into Safe Mode within $\le 60.0\text{ s}$, closing the payload aperture shutter, inhibiting thrusters, acquiring Sun-pointing attitude, and activating S-band telemetry while remaining visible \texttt{SATL} nodes maintain uninterrupted Baltic stereoscopic coverage.

At End of Life ($\ge 7\text{ years}$ \abbr{LEO}; $\ge 15\text{ years}$ \abbr{GEO}, \texttt{R-016}) or critical redundancy exhaustion, every spacecraft executes mandatory active disposal and passivation in compliance with \textbf{\texttt{R-018}} and the \textbf{ESA Zero Debris Policy}. Each \texttt{SATL} executes retrograde perigee-lowering burns using its reserved $\Delta v_{\text{EOL}} \ge 230.0\text{ m/s}$ allocation to ensure complete atmospheric demise within $< 5.0\text{ years}$ ($\ge 99.0\%$ reliability, Design-for-Demise casualty risk $< 10^{-4}$), while each \texttt{SATG} re-orbits to a super-synchronous graveyard orbit $\Delta H \ge +300.0\text{ km}$ above \abbr{GEO} ($\Delta v_{\text{EOL}} \ge 15.0\text{ m/s}$), followed immediately by irreversible venting of residual fluids, battery disconnection, momentum spin-down, and permanent shutdown of all \abbr{RF} and optical transmitters (justifications documented in \textbf{RD-02}).

Table~\ref{tab:sseg_life_profile} summarizes the Space Segment life-profile phases, operational modes, duration envelopes, and governing requirements.

\begin{table}[H]
\centering
\caption{Summary of EWAS Space Segment Life Profile Phases, Operational Modes, and Key Constraints}
\label{tab:sseg_life_profile}
{\scriptsize\setlength{\tabcolsep}{4pt}\renewcommand{\arraystretch}{1.04}%
\arrayrulecolor{TableBorderBlue}%
\rowcolors{2}{white}{TableRowLightBlue}%
\noindent
\begin{tabularx}{\textwidth}{|>{\raggedright\arraybackslash}p{2.5cm}|>{\raggedright\arraybackslash}p{2.0cm}|>{\raggedright\arraybackslash}p{2.7cm}|>{\raggedright\arraybackslash}X|}
    \hline
    \rowcolor{TableHeaderBlue}
    \thead{Lifecycle Phase / Mode} & \thead{Duration / Timing} & \thead{Applicable CI \& Orbit} & \thead{Space Segment Configuration, Key Activities \& Governing Requirements} \\ \hline
    \textbf{Phase 1: On-Ground \abbr{AIT}, Storage \& Transport} & 18--36~months ($\le 24$~mo storage) & \texttt{EM}, \texttt{STM}, \texttt{QM/PFM}, \texttt{FM} (\texttt{SATL} 72x, \texttt{SATG} 6x) & ISO~5 ($\mathrm{GN_2}$ purge) for \abbr{PLD}/\abbr{OISL} optics; ISO~8 for bus; \abbr{HIL} testing (\texttt{R-039}); \abbr{TVAC}, vibration, shock, and EMC; transport loads $\le 3.0\,g$. \\ \hline
    \textbf{Phase 2: Pre-Launch \& Launch Ascent} & Launch campaign + ascent ($\sim 1\text{--}6\text{ h}$) & 8 plane-launches (9~\texttt{SATL}/launch); \texttt{SATG} \abbr{GTO}/\abbr{GEO} & 9-sat dispenser stack ($\le 2{,}800\text{ kg}$); \texttt{SATG} ($\le 3{,}100\text{ kg}$); \abbr{HDRM}s \& aperture covers locked; $\ge 3$ safety inhibits. \\ \hline
    \textbf{Phase 3: \abbr{LEOP} \& In-Orbit Commissioning} & Weeks 1--6 (\texttt{SATL}) \newline Weeks 1--12 (\texttt{SATG}) & Insertion to 1,000~km ($80^\circ$) \& 35,786~km ($0^\circ$) & Bus detumbling ($\le 2\text{--}3^\circ/\text{s}$ to $< 0.1^\circ/\text{s}$); array deployment; 72--96~h bake-out ($310\text{--}320\text{ K}$); cooler cooldown ($65\text{--}95\text{ K}$); orbit insertion \& Walker Delta phasing ($\Delta u_{\text{inter}} = 22.5^\circ$); \abbr{OISL} mesh validation (\texttt{R-040}). \\ \hline
    \textbf{Phase 4 --- Mode N1: Boost Surveillance} & Continuous 24/7 ($\ge 7\text{ yr}$ \abbr{LEO} / $\ge 15\text{ yr}$ \abbr{GEO}) & \texttt{SATG} primary (6x); \texttt{SATL} high-lat.\ suppl. & Persistent $\ge 10\text{ Hz}$ \abbr{WFOV} \abbr{SWIR}/\abbr{MWIR} staring ($P_d > 95\%$, \texttt{R-009}, \texttt{R-010}); onboard proton despiking \& clutter filtering (\texttt{R-019}, \texttt{R-022}); encrypted \emph{GEO Event Report} downlink in $\le 10.0\text{ s}$ (\texttt{R-042}). \\ \hline
    \textbf{Phase 4 --- Mode N2: Agile Stereo Tracking} & Event-driven ($\le 25\text{--}30\text{ s}$ active window) & 2 cued \texttt{SATL} nodes at 1,000~km + \abbr{OISL} & Triggered by Ground \emph{LEO Tracking Request} ($\le 5\text{ s}$ selection/\abbr{OISL} cueing, \texttt{R-043}, \texttt{R-044}); independent gimbal+\abbr{FSM} lock in $< 5.0\text{ s}$ with $< 2\text{ arcsec/s}$ stability (\texttt{R-024}); \abbr{OISL} \abbr{LOS} vector exchange for \emph{Onboard Preliminary Trajectory} (\texttt{R-023}); dual-product downlink (\texttt{R-011}, \texttt{R-012}, \texttt{R-045}). \\ \hline
    \textbf{Phase 4 --- Mode N3: Mandatory \abbr{OISL} Relay} & Continuous 24/7 background mesh & All 72~\texttt{SATL} \& 6~\texttt{SATG} nodes & 4~laser \abbr{OISL} links/\texttt{SATL} (1,550~nm, $\ge 500$~Mbps, $\ge 6{,}500$~km) + \abbr{GEO}--\abbr{LEO} optical trunks ($\ge 1.5\text{--}2.0$~Gbps, 45,000~km); \abbr{QoS} routing; \abbr{COMSEC} (\texttt{R-031}) \& anti-jamming (\texttt{R-032}, \texttt{R-033}). \\ \hline
    \textbf{Phase 4 --- Mode N4: Station-Keeping \& \abbr{CAM}} & Periodic burns (staggered, 1 node at a time) & 1,000~km \abbr{LEO} ($\pm 0.5^\circ$) \& \abbr{GEO} ($\pm 0.05^\circ$) & Orbit maintenance, momentum dumping, and debris Collision Avoidance Manoeuvres ($P_c > 10^{-4}$, \texttt{IF-EXT-04}); staggered to preserve 24/7 Baltic coverage (\texttt{R-008}, \texttt{R-015}). \\ \hline
    \textbf{Phase 4 --- Mode N5: In-Orbit Calibration} & Scheduled maintenance windows & \texttt{SATL} \& \texttt{SATG} fleet & Radiometric blackbody/deep-space calibration, thermo-elastic boresight updates, and secure uplink of archive-refined algorithms (\texttt{R-004}, \texttt{R-041}). \\ \hline
    \textbf{Phase 5: Contingency (\abbr{FDIR} \& Safe Mode)} & L1 $< 100\text{ ms}$; L2 $< 2\text{ s}$; L3 $\le 60\text{ s}$ entry & Any faulted \texttt{SATL} or \texttt{SATG} node & L1 \abbr{EDAC}/\abbr{TMR}/proton despiking/\abbr{SEL} reset; L2 autonomous A/B redundant unit failover (\texttt{R-017}); L3 Safe Mode closes payload shutter, points arrays to Sun, and activates S-band beacon. \\ \hline
    \textbf{Phase 6: \abbr{EOL} Disposal \& Passivation} & Post-mission ($\le 5\text{ yr}$ decay / immediate \abbr{GEO}) & \texttt{SATL} deorbit; \texttt{SATG} $+300\text{ km}$ graveyard & Active \texttt{SATL} retrograde deorbit ($\Delta v_{\text{EOL}} \ge 230\text{ m/s}$, $\ge 99\%$ reliability); \texttt{SATG} graveyard ($\Delta H \ge +300\text{ km}$); full fluid/battery/RF passivation (\texttt{R-018}). \\ \hline
\end{tabularx}}
\end{table}
